\newcommand{\figuraMarcoPeliculaMovil}{%
\begin{figure}[H]
\centering
\begin{tikzpicture}[>=Latex,font=\small,line cap=round,line join=round]
  \fill[cyan!12] (2,1) rectangle (8,4.8);
  \draw[very thick] (2,5.2)--(2,0.6)--(9.4,0.6)--(9.4,5.2);
  \draw[very thick,red!75!black] (8,0.65)--(8,5.15);
  \draw[densely dashed] (9.0,0.65)--(9.0,5.15);
  \draw[->,ultra thick,purple!75!black] (8,2.9)--(10.4,2.9)
    node[right] {$F=\alpha L$};
  \draw[<->,thick] (8,5.55)--(9,5.55) node[midway,above] {$ds$};
  \draw[<->,thick] (7.55,0.65)--(7.55,5.15) node[midway,left] {$L$};
  \node[cyan!50!black] at (4.8,2.9) {una interfaz líquida};
  \node at (5.4,0.15) {$dA=L\,ds,\qquad dW=F\,ds=\alpha\,dA$};
\end{tikzpicture}
\caption{Interfaz con un lado móvil. Al desplazar la barra una distancia $ds$,
el área aumenta en $Lds$ y la tensión superficial de una interfaz ejerce una fuerza
$F=\alpha L$.}
\label{fig:marco-pelicula-movil}
\end{figure}%
}